% !TITLE 水调歌头·明月几时有
% !AUTHOR 苏轼
% !DYNASTY 两宋
% !GENRE 词
% !ORDER 6

\begin{poem}
明月几时有？把酒问青天。
不知天上宫阙\textsuperscript{1}，今夕是何年。
我欲乘风归去，又恐琼楼玉宇\textsuperscript{2}，高处不胜寒\textsuperscript{3}。
起舞弄清影，何似在人间\textsuperscript{4}。

转朱阁\textsuperscript{5}，低绮户\textsuperscript{6}，照无眠。
不应有恨，何事长向别时圆？
人有悲欢离合，月有阴晴圆缺，此事古难全。
但愿人长久，千里共婵娟\textsuperscript{7}。
\end{poem}

\begin{annotations}
\notetext{1}{宫阙：宫殿。阙（què），宫门前的高台建筑。天上的宫殿指月宫。}
\notetext{2}{琼楼玉宇：美玉砌成的楼阁。琼（qióng），美玉；宇，房屋。指想象中的月宫。}
\notetext{3}{高处不胜寒：太高的地方承受不住寒冷。胜（shēng），承受。这句话既写实（月亮确实冷），又暗喻朝堂之上的人情冷暖——苏轼因为政治斗争被贬出京，对“高处”深有体会。}
\notetext{4}{何似在人间：哪里比得上在人间呢。月宫虽美但太冷清，不如人间有烟火气。}
\notetext{5}{朱阁：朱红色的楼阁。月光照在红色的楼阁上。}
\notetext{6}{绮户：雕花的门窗。绮（qǐ），有花纹的丝织品。月光低低地照进雕花窗户。}
\notetext{7}{婵娟：美好的样子，这里指月亮。只要活得长久，哪怕相隔千里，也能共享同一轮明月。这句话成了中国人中秋送祝福最常引用的名句。}
\end{annotations}

\begin{appreciation}
这是中国人过中秋绕不开的一首词，也是写月亮写得最好的词之一。那年中秋，苏轼在密州，喝了一夜的酒，大醉，写这首词寄给弟弟苏辙——题下"兼怀子由"四个字，子由就是苏辙。哥哥写给弟弟，这本书是哥哥写给妹妹，我读到这四个字，觉得千年前有人替这本书写了序。

"明月几时有？把酒问青天。"开篇就问。书里唐人的《把酒问月》问过"青天有月来几时？我今停杯一问之"，李白已经问过一次了，苏轼还要再问一遍，问得更急："明月几时有？"——这月亮到底是从什么时候有的？一个诗人举着酒杯问天，是中国文学里最有名的一次发问之一。

"我欲乘风归去，又恐琼楼玉宇，高处不胜寒。"他想回天上，又怕天上太冷。这里藏着心事：有人说他想回的是朝廷——他离开京城这几年，"高处"的冷是领教过的。"起舞弄清影，何似在人间"：一个人在月光下跳舞，想想还是人间好。人间有酒，有肉，有弟弟。

下阕写月光。"转朱阁，低绮户，照无眠"，月光转过红色的楼阁，低低地照进雕花的窗，照着这个睡不着的人。"不应有恨，何事长向别时圆？"月亮啊月亮，你不该有恨，可为什么总在亲人离别的时候圆？问完，他自己回答："人有悲欢离合，月有阴晴圆缺，此事古难全。"月亮有圆缺，就像人有聚散，自古以来就没法两全。既然没法两全，那就"但愿人长久，千里共婵娟"：都好好活着，隔着千里，共看同一轮月亮。

书里《天问》的屈原追着天问了三百多个问题，天没有回答。一千多年后，苏轼也问天，问完自己给出了答案——不是天告诉他的，是他自己想通的。《生年不满百》说"人生不满百，常怀千岁忧"，苏轼把那个"忧"放下了。从屈原到苏轼，中国人对天的发问，终于等到了一个温柔的答案：求不了全的事，就不求全。

"但愿人长久，千里共婵娟。"这一句苏轼写给弟弟，我厚着脸皮再送给你一回。以后不管你在哪个城市读书，中秋都记得打个电话回家。你抬头看的月亮，和我们看的，是同一轮。
\end{appreciation}
